\documentclass[11pt,a4paper]{article}
\usepackage[margin=2.4cm]{geometry}
\usepackage{amsmath,amssymb,booktabs,array,xcolor}
\usepackage[colorlinks=true,linkcolor=blue!60!black,urlcolor=blue!60!black]{hyperref}
\newcommand{\Lc}{\textsf{\textbf{(c)}}\ }
\newcommand{\Ld}{\textsf{\textbf{(d)}}\ }
\newcommand{\scr}{\texttt{audit\_\allowbreak frames\_\allowbreak defects.py}}
\emergencystretch=2em
\title{Independent audit: defect spectrum, fusion assumptions,\\and the Toda frame}
\author{Quantum Boussinesq solitons project}
\date{4 November 2026}
\begin{document}
\maketitle
\begin{abstract}
The proposed Toda expansion can proceed independently, but the quantum
identification still needs work. \Ld\ An exact symbolic calculation finds a
factor of sixteen missing from the displayed defect sewing determinant.
Its zeros are unaffected, but the stated growing-mode formula also misses a
change of admissible branch. \Lc\ The earlier warning that a KP triangle
passes through unstable defect vacua remains relevant to the claimed quantum
fusion selection rule. The edge-reference subtraction, being independent of
the second soliton, cannot by itself remove the linear ratio dependence
attributed to it in Phase~2h. \Lc
\end{abstract}

\section{Scope and reproducibility}
This note supplements \texttt{independent\_phase2h\_audit.tex} and the separate
charge and determinant audits. Labels (a)--(d) retain their repository meaning;
(c) denotes a project derivation or check and (d) an interpretation or research
judgment. No external literature claim is introduced here. We retain the
good-equation and Corrigan--Zambon dictionaries in \texttt{AGENTS.md}; the
sewing matrix below uses the existing CZ field normalisation. The audit does
not choose a regularisation or quantum phase convention.

All displayed new mathematical results are checked in
\texttt{phase2/scripts/}\scr, with exact SymPy identities and 70-digit mpmath
calculations. Its JSON companion records numerical results and SHA-256 hashes
of the audited inputs. The previously existing
\texttt{oneloop\_defect\_scattering\_and\_stability.py} also completes with
\texttt{ALL CHECKS PASSED}; its tests do not expose the issues below.

\section{The defect determinant: check away from its zeros}
\paragraph{Corrected normalisation \Lc.}
Let $s=\eta\sigma$, and let $\omega=i g$ with $g>0$ be a purely growing
frequency. The decaying half-line spatial exponents satisfy
\[
 \kappa_1^2+\kappa_2^2=1,\qquad
 \kappa_1\kappa_2=g,\qquad
 S=\kappa_1+\kappa_2=\sqrt{1+2g},\qquad
 d=\kappa_1-\kappa_2,\quad d^2=2-S^2.
\]
For the actual $4\times4$ linearised sewing matrix implemented in Phase~2g,
at $\eta=+1$, the exact factorisation is
\begin{equation}\label{eq:det}
 \frac{\det M}{(\kappa_1-\kappa_2)^2}
 =-\frac{2}{3}S^2(S^2-1)P_-(S;s)P_+(S;s),
\end{equation}
where the reported polynomials themselves are correct:
\[
 P_\mp(S;s)=S^2-\frac{\sqrt3\mp1}{2}sS
                   +\frac{s^2-2\mp\sqrt3}{2}.
\]
The coefficient in Phase~2g, Eq.~(\texttt{det\_factorisation}), is
$-1/24$, sixteen times too small. The script proves \eqref{eq:det}
symbolically by replacing each half-line pair of columns by its sum and
its difference divided by $d$, expanding the determinant, and imposing
$d^2=2-S^2$. It also compares an independently evaluated determinant to
the existing matrix at six nonzero samples:
$s=-1.5,0.5,1.9$, $g=0.2,0.7$.
The ratio of actual to displayed determinant is sixteen in every case.

This constant does not change the eigenvalue zeros, and it is not by itself
a computed error in the physical one-loop determinant. \Ld\ The existing
factorisation test evaluates the two sides at a characteristic root, where
both vanish. It therefore does not test the prefactor. \Lc

\paragraph{Admissible roots \Lc.}
For $-2<s<2$, retain only real roots with $S>1$ from
\begin{align*}
 S_{-,\pm}&=\frac{(\sqrt3-1)s\pm(\sqrt3+1)\sqrt{4-s^2}}4,\\
 S_{+,\pm}&=\frac{(\sqrt3+1)s\pm(\sqrt3-1)\sqrt{4-s^2}}4,
 \qquad g=\frac{S^2-1}{2}.
\end{align*}
Squaring an inadmissible root loses the decaying-sheet condition. The
two positive growth rates at $s=1.9$ are approximately
\[
 g=0.200256\ldots,\qquad 0.496894\ldots;
\]
the JSON gives all digits. The smaller mode comes from $S_{+,-}>1$;
$S_{-,+}<1$ has ceased to describe a growing mode. The displayed
frequency pair in Phase~2g instead includes the negative of that smaller
rate. Direct substitution of both admissible rates into the sewing
determinant gives zero to the 70-digit calculation's tolerance.

At $s=2$ one has
\[
 P_+(S;2)=\left(S-\frac{\sqrt3+1}{2}\right)^2,
 \qquad g=\frac{\sqrt3}{4}.
\]
Thus the admissible roots coalesce at the same positive growth rate.
The displayed frequency pair gives opposite signs there, contradicting
the coalescence stated immediately afterwards. The existing script tests
$s=1.5$ and then uses a separate expression at $s=2.5$; it skips the
intervening branch change. These are explicit counterexamples to that
formula as a complete spectrum, not a replacement completeness proof.

\section{Fusion: an algebraic identity does not supply missing states}
\paragraph{Stable-vertex check \Lc.}
The KP triple may be parametrised as
\[
 q_j=\frac{1}{\sqrt3}\cos\left(\vartheta+\frac{2\pi j}{3}\right),
 \qquad \sigma_j=-2\sqrt3 q_j,\qquad j=0,1,2.
\]
The script verifies the pair relation
$q_i^2+q_iq_j+q_j^2=1/4$ and counts admissible growing modes on
198 oriented, nondegenerate sampled triples. Each has exactly one vertex
with no growing mode, agreeing with the previously established condition
$\eta\sigma<-1$. This is a sample check, not a new proof of spectral
completeness at thresholds.

\texttt{defect\_semiclassics.tex}, its stability section, explicitly
downgraded the proposed exact quantum triangle bootstrap because its two
excited defect vacua are unstable. Phase~2g proves a continuum phase
coboundary identity and then calls it a decisive quantum selection rule;
that step does not restore stable intermediate states. \Ld\ Nor does a
continuum identity determine the contribution of any missing discrete modes
or the normalisation of the quantum defect transitions. An analytic
continuation involving resonances might still obey useful identities, but
that needs its own derivation. We retain the algebraic fusion result and
do not treat it alone as a physical exclusion of a Gamma factor or a
selection of an exchange-phase constant. \Ld

\section{The edge subtraction and the quantum frame}
\paragraph{What the code subtracts \Lc.}
In \texttt{toda\_edge\_twoterm.py}, the phase numerator is
\[
 \Phi(p_1,p_2)-\Phi(p_1,e_2)
 -(-2\kappa_1)\left(\frac{P_{n2}}a-\Pi_2\right).
\]
The frame term uses the classical quantities
$P_{n2}=4b(\kappa_2\cosh\kappa_2-\sinh\kappa_2)$ and
$\Pi_2=2b\sinh\kappa_2$ in the right-moving sector. With
$\kappa_j=\sqrt{y_j\hbar_T/b}$, the script verifies
\[
 \lim_{\hbar_T\to0}
 \frac{(-2\kappa_1)(P_{n2}/a-\Pi_2)}{\hbar_T}
 =4\sqrt{y_1y_2}.
\]
This establishes the implemented classical subtraction. It does not
establish that the same transformation has no quantum correction.
\texttt{toda\_edge.tex} explicitly retains that caveat, as does
\texttt{matching\_toda.py}. A complete comparison to continuum quantum
scattering still needs a common quantum definition of the frame and of
the physical charge labels. \Ld

\paragraph{An edge value cannot subtract a slope \Lc.}
Phase~2h attributes the $+4r/\pi$ term in its candidate limiting residual
to subtracting the edge reference $-\Phi(p_1,e_2)/\hbar_T$, saying that
this removes the linear small-$r$ slope. At fixed $y_1=Y$ the ground
state, $p_1$, and $e_2$ are independent of $y_2=r^2Y$. Therefore this
subtraction is a function $C(Y)$ alone, and exactly
\[
 \partial_r\{F(Y,r)-C(Y)\}=\partial_r F(Y,r).
\]
Equivalently, differences between any two ratios are unchanged. It cannot
alone add the nonconstant function $4r/\pi$. This invalidates that
explanation of the term, without disproving the candidate final residual:
the term might instead come from the dressed phase, charge conversion,
or another part of a derivation not yet supplied. \Ld\ The earlier audit
already records that the script defines the purported loop asymptotic
by subtracting the classical contribution from the proposed target.

\section{Checks, failures recorded, and verdict}
\begin{center}\small
\begin{tabular}{p{.57\textwidth}p{.34\textwidth}}\toprule
Claim or scope & Verification\\\midrule
Corrected sewing determinant, away from zeros (c) & Exact polynomial identity;
six 70-digit samples\\
Growing-mode counterexample and coalescence (c) & Root filtering, direct
determinants, exact double root\\
One stable vertex on sampled triples (c) & 198 oriented nondegenerate triples\\
Classical frame limit; edge subtraction preserves slope (c) & Exact SymPy identities\\
Physical quantum frame and fusion selection (d) & Still require derivation\\\bottomrule
\end{tabular}
\end{center}

\paragraph{Failures recorded.}
\begin{enumerate}
\item The displayed factorisation fails away from its roots by a factor
of sixteen. The stated frequency pair fails the admissible growing-branch
test at $s=1.9$ and the stated coalescence at $s=2$.
\item An exploratory call to mpmath's LU determinant raised a
\texttt{TypeError} at an exactly singular matrix. The independent checker
uses the exact 24-term determinant definition, retaining the same matrix
and tolerance; no eigenvalue condition is dropped.
\item A direct symbolic determinant computation was interrupted after
stalling. Symmetric and antisymmetric column combinations reduce its size
and prove the same off-shell identity exactly.
\end{enumerate}

\end{document}
